%======================================================================
%  Hardware platform: SD650 V3 and Intel Xeon CPU Max
%======================================================================

\section{Xeon Max and the SD650 V3}
\label{sec:xeon-max-platform}

\subsection{The unit of analysis}

The Lenovo ThinkSystem SD650~V3 is a dense, direct-water-cooled server
tray rather than a single compute node.  One tray contains two
independent dual-socket nodes, and six trays fit in a 6U DW612S
enclosure~\cite{lp1603}.  This distinction matters when reporting
quantities: the CRESCO8 allocation is a \emph{node}; memory locality is
defined first at the \emph{socket}; and the scheduler may allocate
individual cores within the node.  \tabref{tab:hardware-hierarchy} sets
out the four levels and what each one is good for.

\begin{table}[htbp]
\centering
\small
\tabsetup
\caption{Hardware hierarchy and the properties relevant to this study.
Product capability is separated from the installed CRESCO8
configuration described in \secref{sec:cresco8}
~\cite{lp1603,cresco8doc}.}
\label{tab:hardware-hierarchy}
\begin{tabularx}{\textwidth}{@{}lYY@{}}
\toprule
\headrow
\textbf{Level} & \textbf{Composition} & \textbf{Relevance} \\
\midrule
DW612S enclosure
  & Six SD650~V3 trays (12 independent nodes)
  & Shared power and cooling infrastructure; not an allocation unit. \\
SD650~V3 tray
  & Two independent dual-socket nodes
  & High-density packaging; the two nodes do not form one shared-memory
    system. \\
Compute node
  & Two Xeon processors connected by Intel UPI
  & Scheduler allocation and the largest cache-coherent NUMA system used
    by the single-node experiments. \\
Xeon Max socket
  & CPU cores, 64\,GB on-package HBM2e and eight DDR5 channels
  & Primary locality domain and the cleanest unit for comparing HBM with
    DDR5. \\
\bottomrule
\end{tabularx}
\end{table}

The product guide covers several processor, memory, storage and network
options.  Most of those options are irrelevant to the thesis.  The
important configuration is the Xeon CPU Max 9480 installed on CRESCO8:
56 physical cores per socket, 64\,GB of HBM2e, eight DDR5-4800 channels,
and four UPI links~\cite{intel9480,lp1603}.  A two-socket node therefore
provides 112 physical cores and 128\,GB of nominal HBM capacity.  The
node also contains DDR5, but its installed capacity and channel
population are properties of CRESCO8, not of the product family.

\begin{figure}[htbp]
\centering
\begin{tikzpicture}[
  font=\small,
  socket/.style={draw=black!65, fill=neutralcol, rounded corners=3pt,
                 minimum width=4.5cm, minimum height=1.35cm,
                 align=center},
  hbm/.style={draw=hbmcol!85!black, fill=hbmcol!18, rounded corners=2pt,
              minimum width=4.5cm, minimum height=0.85cm, align=center},
  ddr/.style={draw=ddrcol!85!black, fill=ddrcol!14, rounded corners=2pt,
              minimum width=4.5cm, minimum height=0.85cm, align=center},
  link/.style={{Stealth}-{Stealth}, thick, black!65}
]
  \node[socket] (s0) at (-3.8,0)
    {\textbf{Socket 0}\\56 cores, AMX, LLC};
  \node[socket] (s1) at (3.8,0)
    {\textbf{Socket 1}\\56 cores, AMX, LLC};
  \node[hbm, below=0.55cm of s0] (h0)
    {64\,GB HBM2e\\high bandwidth};
  \node[hbm, below=0.55cm of s1] (h1)
    {64\,GB HBM2e\\high bandwidth};
  \node[ddr, above=0.55cm of s0] (d0)
    {8 DDR5 channels\\high capacity};
  \node[ddr, above=0.55cm of s1] (d1)
    {8 DDR5 channels\\high capacity};

  \draw[link] (s0) -- node[above, fill=white, inner sep=2pt]{UPI} (s1);
  \draw[link] (s0) -- (h0);
  \draw[link] (s1) -- (h1);
  \draw[link] (s0) -- (d0);
  \draw[link] (s1) -- (d1);

  \node[draw=black!35, dashed, rounded corners=4pt, inner sep=9pt,
        fit=(d0)(d1)(h0)(h1),
        label={[black!60]below:{one cache-coherent compute node}}] {};
\end{tikzpicture}
\caption{Conceptual locality map of a CRESCO8 Xeon Max node.  Each
socket has local HBM and DDR5; access to the other socket traverses UPI.
The diagram represents locality, not physical package routing.
Author's synthesis based on the platform and cluster documentation
~\cite{lp1603,cresco8doc}.}
\label{fig:node-locality}
\end{figure}


\subsection{What the hardware looks like}
\label{sec:hardware-photographs}

Three photographs from the product documentation make
\tabref{tab:hardware-hierarchy} concrete.  They are worth a moment,
because the packaging explains several constraints that reappear later
as experimental limitations.

\begin{figure}[htbp]
\centering
\includegraphics[width=0.94\textwidth]{sd650v3-front-view.png}
\caption{Front of one SD650~V3 tray.  The two servers are independent
computers that happen to share a chassis: each has its own network
ports, its own operator panel and its own management interface, and
they share only power and water.  What is \emph{not} present is as
informative as what is.  There is no video output, no removable media
and no local console, because a node of this kind is never approached
directly.  Reproduced from the Lenovo ThinkSystem SD650~V3 product
guide~\cite{lp1603}.}
\label{fig:sd650-front}
\end{figure}

\figref{fig:sd650-front} shows the tray from the front.  Two servers,
side by side, each with a pair of 25\,GbE ports for the high-speed
fabric and a 1\,GbE port for management.  The experiments in this thesis
touch none of these interfaces: work arrives through the scheduler and
results leave through the file system.

\begin{figure}[htbp]
\centering
\includegraphics[width=0.94\textwidth]{sd650v3-internal-view.png}
\caption{Inside the tray.  Each of the two servers carries two
processors, and each processor is flanked by eight memory DIMMs.  The
copper cold plates and the loop that joins them replace the fans
entirely.  The high-bandwidth memory that this thesis is about is
nowhere visible: it sits inside the processor package, beneath the cold
plate, which is why its capacity is fixed and why it cannot be
configured, replaced or upgraded.  Reproduced from the Lenovo
ThinkSystem SD650~V3 product guide~\cite{lp1603}.}
\label{fig:sd650-internal}
\end{figure}

\figref{fig:sd650-internal} is the more useful of the two.  The layout
of one server, two processors with their own DIMMs on either side, is
the physical form of the locality argument made throughout this section:
a thread running on the left-hand processor reaches the DIMMs beside it
quickly and the DIMMs on the right slowly, and no amount of software can
change the distance.  The absence of fans is also visible, and it
matters for measurement rather than for acoustics, since a water-cooled
socket holds its clock under sustained load instead of ramping fans and
throttling.

\begin{figure}[htbp]
\centering
\includegraphics[width=0.80\textwidth]{dw612s-enclosure-nodes.png}
\caption{Six trays in one 6U DW612S enclosure, giving twelve independent
two-socket nodes.  A scheduler allocation yields one of these, not the
enclosure and not the rack.  Reproduced from the Lenovo ThinkSystem
SD650~V3 product guide~\cite{lp1603}.}
\label{fig:dw612s-enclosure}
\end{figure}

\figref{fig:dw612s-enclosure} completes the picture: twelve nodes in one
enclosure, and a few dozen enclosures in the machine room.  A job that
asks for a node receives one of the trays in this photograph, with its
neighbours in the enclosure running unrelated work.  That is the reason
\secref{sec:first-allocation} insists on recording which node was
obtained.

\subsection{Compute capability, and why it is not the binding constraint}

Xeon Max is a Sapphire Rapids processor with conventional x86 cores,
AVX-512, and Intel Advanced Matrix Extensions (AMX).  AMX adds per-core
tile registers and matrix-multiply instructions for BF16 and INT8
operands~\cite{intelmax,na2024llmcpu}, and on well-shaped
matrix--matrix products it is a large win, which is why prompt
processing benefits from it and single-token generation largely does
not.

The detail that matters for a quantized workload is what AMX does not
provide.  There is no native INT4 dot-product instruction, so a 4-bit
GGUF model cannot be fed to the matrix engine in the format it is stored
in.  \figref{fig:amx-path} shows the consequence: every weight must
pass through an unpack and dequantize step before it reaches a tile, and
that step sits precisely between the memory tier and the matrix unit.
The arrangement has an awkward property for this thesis.  Lower-bit
formats reduce the bytes moved, which is what makes them attractive on a
bandwidth-limited machine, but they add work on the path from memory to
the matrix engine, and they push the operand format further from
anything the engine can consume directly.  Whether a given model
benefits from AMX at all therefore depends on the runtime, the storage
format, the packing scheme and the matrix dimensions, none of which are
properties of the processor.

\begin{figure}[htbp]
\centering
\begin{tikzpicture}[
  font=\small,
  mem/.style={draw=hbmcol!85!black, fill=hbmcol!14, rounded corners=3pt,
              minimum width=2.3cm, minimum height=1.5cm, align=center},
  stage/.style={draw=black!65, fill=neutralcol, rounded corners=3pt,
                minimum width=2.7cm, minimum height=1.5cm, align=center},
  eng/.style={draw=black!70, fill=white, rounded corners=3pt,
              minimum width=2.4cm, minimum height=1.5cm, align=center},
  flow/.style={-{Stealth}, thick, black!70},
  edgelbl/.style={font=\scriptsize, align=center, text=black!70}
]
  \node[mem]   (m) at (0,0)    {HBM2e\\\emph{or}\\DDR5};
  \node[stage] (u) at (4.3,0)  {unpack \&\\dequantize\\{\footnotesize AVX-512}};
  \node[stage] (t) at (8.6,0)  {tile registers\\{\footnotesize BF16 / INT8}};
  \node[eng]   (x) at (12.1,0) {AMX\\\textsc{tmul}};

  \draw[flow] (m) -- node[above=1pt, edgelbl] {4-bit}
                    node[below=1pt, edgelbl] {packed} (u);
  \draw[flow] (u) -- node[above=1pt, edgelbl] {widened} (t);
  \draw[flow] (t) -- (x);

  \draw[dashed, thick, black!45, -{Stealth}]
    (m.north) .. controls +(2.4,1.6) and +(-2.8,1.6) .. (x.north);
  \node[edgelbl, text=black!55, fill=white, inner sep=2pt] at (6.05,1.62)
    {no native INT4 path};
\end{tikzpicture}
\caption{Path taken by a quantized weight from a memory tier to the
matrix engine on Xeon Max.  The unpack stage is unavoidable for 4-bit
formats and lies between the two things this thesis compares.  At batch
one the right-hand operand is a vector, so the tile is largely empty and
the engine runs far below its rated throughput; the step is paced
instead by the arrival of weights from the left.  Author's synthesis
from the Xeon Max product documentation and the AMX description in Na et
al.~\cite{intelmax,na2024llmcpu}.}
\label{fig:amx-path}
\end{figure}

\textbf{Peak instruction throughput is not an application result.}  Small
matrices, batch-one matrix--vector operations, packing overhead and
memory stalls can all leave the matrix units idle, and the advertised
maximum turbo frequency is a single-core limit rather than the frequency
of a socket running sustained AMX kernels.  Frequency, package power and
throttling state have to be observed during the experiments rather than
read off a specification table.


\subsection{The memory hierarchy}
\label{sec:memory-hierarchy}

HBM and DDR5 trade different resources.  HBM provides much greater
aggregate bandwidth and 64\,GB of nominal capacity per socket.  DDR5
provides substantially more capacity and can have lower unloaded access
latency.  HBM is therefore not simply ``faster memory''.  It is most valuable
under two conditions that have to hold \emph{together}: the workload
issues enough concurrent requests to approach its bandwidth, \emph{and}
the bandwidth-critical working set is placed locally.  Either one alone
is not enough.

Published STREAM results illustrate both the advantage and the need for
local measurement.  Na et al.\ report 588\,GB/s for HBM and
233.8\,GB/s for DDR5 on a Xeon Max 9468 socket
~\cite{na2024llmcpu}.  Ibeid et al.\ report 680--840\,GB/s for HBM and
235--245\,GB/s for DDR5 on a related Xeon Max 9470C
~\cite{ibeid2025aurora}.  These sustained values are 2.5--3.4 times
apart, but the HBM results themselves differ by more than 15\%.  SKU,
clustering mode, thread placement and benchmark configuration are
plausible causes.  None of these numbers should be used as the bandwidth
of CRESCO8 before the Max 9480 nodes are measured.

Three paths must be distinguished in every result:

\begin{itemize}
  \item \textbf{local HBM or local DDR5}, reached from cores on the same
        socket;
  \item \textbf{remote memory}, reached through UPI from the other
        socket; and
  \item \textbf{mixed placement}, in which pages, threads or model
        objects span several locality domains.
\end{itemize}

\figref{fig:node-locality} draws these three paths.  An apparently
node-wide 128\,GB HBM allocation is consequently two 64\,GB local
tiers.  A workload uses the aggregate capacity and
bandwidth only if data and computation are balanced across sockets.

\begin{table}[htbp]
\centering
\small
\tabsetup
\caption{Storage levels visible to one Max 9480 socket, ordered by the
question that matters here: who decides what occupies them.  Only the
last three rows are under software control, and only in Flat mode.
Cache capacities follow the SKU specification and should be confirmed
with \texttt{lscpu} on the allocated node~\cite{intel9480,lp1603}.}
\label{tab:storage-levels}
\begin{tabularx}{\textwidth}{@{}P{2.35cm}P{3.05cm}P{3.15cm}Y@{}}
\toprule
\headrow
\textbf{Level} & \textbf{Capacity} & \textbf{Placement decided by} &
\textbf{Why it matters here} \\
\midrule
L1 data & 48\,KB per core & Hardware & Too small to hold any part of a
  weight tensor. \\
L2 & 2\,MB per core & Hardware & Holds tiles and activations; invisible
  to a placement policy. \\
Last-level cache & 112.5\,MB, shared & Hardware replacement policy &
  The tier Zhang et al.\ target by rebuilding the
  runtime~\cite{zhang2026cacheresident}; not addressable directly. \\
\addlinespace[2pt]
HBM2e & 64\,GB, four stacks & \textbf{Software}, in Flat mode &
  The bandwidth tier, and the object of this study. \\
DDR5 & Eight channels; installed capacity is a CRESCO8 property &
  \textbf{Software}, in Flat mode & The capacity tier and the control
  condition. \\
Remote memory & The other socket's tiers & \textbf{Software}, through
  thread and page binding & Reached over UPI; a placement failure rather
  than a design choice. \\
\bottomrule
\end{tabularx}
\end{table}

\tabref{tab:storage-levels} also locates this thesis against the
closest competing work.
Zhang et al.\ obtain large gains by keeping weights resident in the
last-level cache~\cite{zhang2026cacheresident}, one row above the tier
studied here, but that row is not addressable, so their result required
rebuilding the runtime around it.  The HBM row can be
addressed directly, which is what makes the comparison in this thesis
possible without touching the runtime at all.

\subsection{Bandwidth is only reachable with enough concurrency}
\label{sec:littles-law}

A bandwidth number describes what the memory system can deliver when it
is kept busy, and nothing about whether a given program manages to keep
it busy.  The relationship is Little's law: the number of bytes a
program must have in flight to sustain a bandwidth \(B\) against a
latency \(L\) is

\begin{equation}
\label{eq:little}
  C = B \times L .
\end{equation}
\eqcaption{Little's law applied to a memory system.  It converts a bandwidth figure into a requirement on the program: this many bytes must be outstanding at every instant, or the bandwidth is unreachable no matter which tier holds the data.}

Substituting the published Xeon Max figures into \eqnref{eq:little}
makes the difference between the two tiers concrete.  For HBM at 840\,GB/s and 134.9\,ns, roughly
113\,KB must be outstanding at all times, which is about 1{,}770 cache
lines, or close to 32 lines per core across 56 cores.  For DDR5 at 245\,GB/s
and 112.9\,ns the requirement is roughly 28\,KB, about 430 lines, or
under 8 per core~\cite{ibeid2025aurora}.  The faster tier does not merely
offer \emph{more} bandwidth; it demands roughly \textbf{four times the
memory-level parallelism} before that bandwidth becomes reachable at
all.

That demand is not obviously satisfiable.  Thirty-two concurrent line
fills per core is the same order of magnitude as the outstanding-miss
resources a single core has available, so approaching HBM peak requires
almost every one of them to stay occupied, helped by hardware
prefetchers that work well on strided streams and poorly on dependent
ones.  The exact resource counts for this microarchitecture must be
confirmed rather than assumed, but the shape of the constraint does not
depend on the precise number.

This is the mechanism behind the most important negative result in the
tiered-memory literature.  Peng et al.\ found that HPC codes with
regular access patterns gained up to 3$\times$ from MCDRAM while codes
with irregular, latency-bound patterns could run \emph{slower} in the
high-bandwidth tier, and that adding hardware threads, which is to say adding concurrency,
recovered much of the loss~\cite{peng2017hybridmemory}.  Low-batch autoregressive decode has
exactly the awkward shape: the weight stream is regular and prefetchable,
but the dependent chain through each layer limits how far ahead the core
can run, and the attention path over a growing KV cache is less regular
still.

Two consequences follow for the experiments.  Thread count is not a
nuisance parameter to be fixed at ``all cores'' but a first-class
factor, because it directly sets available concurrency; and a
disappointing HBM result must be diagnosed by measuring achieved
bandwidth rather than reported as an absence of effect.  A run that
reaches 250\,GB/s on HBM has \emph{not} shown that HBM does not help;
it has shown that \textbf{the program never asked HBM for anything DDR5
could not also have delivered}.

\subsection{HBM memory modes}
\label{sec:hbm-modes}

Xeon Max supports three ways to expose HBM, shown side by side in
\figref{fig:hbm-modes}~\cite{lp1603,intelmaxtuning}.  The mode is
selected in firmware and is normally fixed for the lifetime of a boot.

\begin{figure}[htbp]
\centering
\begin{tikzpicture}[
  font=\small,
  cpu/.style={draw=black!65, fill=neutralcol, rounded corners=2pt,
              minimum width=3.9cm, minimum height=0.75cm, align=center},
  hbm/.style={draw=hbmcol!85!black, fill=hbmcol!18, rounded corners=2pt,
              minimum width=3.9cm, minimum height=0.75cm, align=center},
  ddr/.style={draw=ddrcol!85!black, fill=ddrcol!14, rounded corners=2pt,
              minimum width=3.9cm, minimum height=0.75cm, align=center},
  absent/.style={draw=black!35, dashed, text=black!50, rounded corners=2pt,
                 minimum width=3.9cm, minimum height=0.75cm, align=center},
  link/.style={{Stealth}-{Stealth}, thick},
  note/.style={font=\footnotesize, text width=4.2cm, align=center}
]
  \node[font=\bfseries] at (0,1.15) {HBM-only};
  \node[cpu] (c0) at (0,0) {CPU cores};
  \node[hbm] (h0) at (0,-1.25) {HBM};
  \node[absent] (d0) at (0,-2.5) {DDR5 absent};
  \draw[link] (c0) -- (h0);
  \node[note] at (0,-3.55)
    {One visible memory tier; capacity limited to HBM.};

  \node[font=\bfseries] at (5.2,1.15) {Flat};
  \node[cpu] (c1) at (5.2,0) {CPU cores};
  \node[hbm] (h1) at (5.2,-1.25) {HBM NUMA nodes};
  \node[ddr] (d1) at (5.2,-2.5) {DDR5 NUMA nodes};
  \draw[link] (c1) -- (h1);
  \draw[link] (c1.east) .. controls +(0.65,-0.55) and
                                  +(0.65,0.55) .. (d1.east);
  \node[note] at (5.2,-3.55)
    {Both tiers are addressable; software controls placement.};

  \node[font=\bfseries] at (10.4,1.15) {Cache};
  \node[cpu] (c2) at (10.4,0) {CPU cores};
  \node[hbm] (h2) at (10.4,-1.25) {HBM cache};
  \node[ddr] (d2) at (10.4,-2.5) {DDR5 address space};
  \draw[link] (c2) -- (h2);
  \draw[link] (h2) -- (d2);
  \node[note] at (10.4,-3.55)
    {Hardware-managed, direct-mapped cache in front of DDR5.};
\end{tikzpicture}
\caption{The three Xeon Max memory modes, shown for one socket.  Flat
mode is the only mode that exposes HBM and DDR5 as separately
controllable memory tiers.  Author's synthesis from the Xeon Max product
and tuning guides~\cite{lp1603,intelmaxtuning}.}
\label{fig:hbm-modes}
\end{figure}

\begin{description}
  \item[HBM-only mode.] The system contains no DDR5 and exposes HBM as
  ordinary system memory.  It is simple for software but imposes a hard
  capacity limit.  It is not the documented CRESCO8 configuration.

  \item[Flat mode.] HBM and DDR5 share the physical address space but
  appear as distinct NUMA memory nodes.  Applications can allocate
  directly from a tier, or an external policy can bind pages with
  \texttt{mbind}, libnuma or \texttt{numactl}.  Allocation is not enough:
  first touch, page migration and memory mapping can change actual page
  residency, which must be verified.

  \item[Cache mode.] HBM is a hardware-managed, direct-mapped cache in
  front of DDR5.  Applications see the DDR5 capacity without placing
  objects explicitly, but software cannot choose which objects remain
  in HBM.  Conflict misses, warm state and a working set near the cache
  capacity can make performance workload-dependent
  ~\cite{ibeid2025aurora}.
\end{description}

Flat mode provides the strongest experimental design because local HBM
and local DDR5 can be compared on the same socket without changing the
processor or executable.  A Flat-versus-Cache comparison answers a
different question: it measures explicit placement against a hardware
cache policy.

\subsection{Clustering and NUMA}

Memory mode is independent of the processor's clustering mode.  In
Quadrant mode, one socket is normally presented as one compute locality
domain.  Sub-NUMA Clustering (SNC4) exposes four domains per socket,
each associated with a subset of cores, cache, DDR channels and HBM.
SNC4 can increase bandwidth and reduce local latency for software that
places workers and data per quadrant, but it also makes remote accesses
easier to create.  In Flat+SNC4, the nominal 64\,GB HBM socket is divided
into four 16\,GB locality domains.

This distinction explains why a mode cannot be labelled universally
``best''.  Ibeid et al.\ obtain the highest STREAM bandwidth with SNC4
and explicit locality, whereas Na et al.\ observe worse inference
performance and more remote traffic under SNC4
~\cite{ibeid2025aurora,na2024llmcpu}.  The result depends on the
software's placement policy, not only on the processor setting.

\subsection{Cooling and instrumentation}

Direct water cooling allows two high-power processors to operate in a
dense node, but it does not guarantee a fixed frequency or energy
efficiency.  The relevant quantities are measured socket frequency,
temperature, package power and throttling state.  Product-guide claims
about facility heat recovery or cooling savings are not evidence about
inference performance.

\begin{figure}[htbp]
\centering
\includegraphics[width=0.82\textwidth]{dw612s-water-manifold.png}
\caption{Rear of the enclosure.  Water is distributed to the trays
through manifolds, and even the power supplies are water-cooled, so the
enclosure has no air path at all.  For this thesis the consequence is
narrow but useful: thermal conditions are set by facility water
temperature rather than by neighbouring workloads, which removes one
source of run-to-run variation.  Reproduced from the Lenovo ThinkSystem
SD650~V3 product guide~\cite{lp1603}.}
\label{fig:dw612s-manifold}
\end{figure}

\figref{fig:dw612s-manifold} shows how the water reaches the trays.
The SD650~V3 includes XClarity management and a TI INA226-based
power/energy meter with 100\,Hz sampling capability
~\cite{lp1603}.  This establishes a platform capability, not user
access on CRESCO8.  Intel RAPL, management-controller telemetry and
whole-node power meters cover different system boundaries; they must
not be combined without first validating their scope, units, sampling
and permissions.

\subsection{Consequences for the thesis}

Five requirements follow from the hardware, and each of them exists
because a specific thing can go wrong.

\begin{enumerate}
  \item Establish the active memory and clustering modes from the
        allocated node rather than from the product guide.  Flat mode is
        a precondition for the primary comparison, not a default.
  \item Measure sustained local HBM and DDR5 behaviour before drawing a
        Roofline or interpreting any inference result.  The published
        figures for this processor family differ by more than 15\%
        between studies, which is too much to inherit.
  \item Treat thread count as a factor rather than a constant.  It sets
        the memory-level parallelism available to the run, and therefore
        how much of the HBM bandwidth the workload can reach at all
        (\secref{sec:littles-law}).
  \item Bind threads and pages, then verify residency and remote traffic
        rather than trusting the command line.  A memory policy records
        an intention; only a page count records an outcome.
  \item Record frequency, thermal and power state throughout, so that a
        placement effect cannot be confused with a change in operating
        conditions.
\end{enumerate}

The common thread is that this processor offers several ways to obtain a
result that looks like a memory-tier effect and is not one.  Most of the
methodological weight in Part~III goes on ruling those out.
